% Generated by `taskgen course examples` from dlp.ic_linear_system@singular (seed 2).
% Do not edit: change examples.yaml and regenerate.
\begin{example}[система по модулю $p - 1$, вырожденная по модулю $2$]\label{ex:ic-system-singular}
Пусть $p = 109$, $g = 56$ --- первообразный корень по модулю $p$, факторная база $F = \{2, 3, 5\}$. Найдены соотношения\[g^{19} \equiv 2^{5} \cdot 3,\quad g^{24} \equiv 3 \cdot 5^{2},\quad g^{98} \equiv 2^{2} \cdot 5^{2},\quad g^{104} \equiv 3 \cdot 5\pmod{109}.\]Определяют ли первые 3 из них логарифмы $L_\ell = \log_g \ell$ однозначно? Найти $L_\ell$, $\ell \in F$, используя все соотношения (систему решать по модулю примарных делителей $p - 1$ и собирать по КТО).

\par\noindent\textit{Решение.} Неизвестные $L_\ell$ определены по модулю $p - 1 = 108 = 2^{2} \cdot 3^{3}$. Система:\[\begin{array}{rcl}(1)\colon\ 5L_{2} + L_{3} & \equiv & 19 \\ (2)\colon\ L_{3} + 2L_{5} & \equiv & 24 \\ (3)\colon\ 2L_{2} + 2L_{5} & \equiv & 98 \\ (4)\colon\ L_{3} + L_{5} & \equiv & 104 \\ \end{array}\pmod{108}\]Определитель первых 3 строк: $D = 14$, $\gcd(D, 108) = 2$. По модулю $2$ матрица первых 3 соотношений вырождена: они \emph{не} определяют $L_\ell$ однозначно (решений по модулю $q^e$ несколько), нужно соотношение (4).\par Метод Гаусса по модулю каждого $q^e$; ведущий элемент должен быть обратим, т.\,е. не делиться на $q$ (столбцы --- коэффициенты при $L_{2}$, $L_{3}$, $L_{5}$ и правая часть; над стрелкой --- исключаемая переменная, номер ведущей строки в текущей матрице и обратный к ведущему коэффициенту).\par\noindent По модулю $4$:\[\left(\begin{array}{rrr|r}1 & 1 & 0 & 3 \\ 0 & 1 & 2 & 0 \\ 2 & 0 & 2 & 2 \\ 0 & 1 & 1 & 0 \\ \end{array}\right)\xrightarrow{L_{2}:\ (1),\ 1^{-1} = 1}\left(\begin{array}{rrr|r}1 & 1 & 0 & 3 \\ 0 & 1 & 2 & 0 \\ 0 & 2 & 2 & 0 \\ 0 & 1 & 1 & 0 \\ \end{array}\right)\]\[\xrightarrow{L_{3}:\ (2),\ 1^{-1} = 1}\left(\begin{array}{rrr|r}1 & 0 & 2 & 3 \\ 0 & 1 & 2 & 0 \\ 0 & 0 & 2 & 0 \\ 0 & 0 & 3 & 0 \\ \end{array}\right)\xrightarrow{L_{5}:\ (4),\ 3^{-1} = 3}\left(\begin{array}{rrr|r}1 & 0 & 0 & 3 \\ 0 & 1 & 0 & 0 \\ 0 & 0 & 1 & 0 \\ 0 & 0 & 0 & 0 \\ \end{array}\right)\]$\Rightarrow$ $L_{2} \equiv 3$, $L_{3} \equiv 0$, $L_{5} \equiv 0$ $(\textrm{mod}\ 4)$.\par\noindent По модулю $27$:\[\left(\begin{array}{rrr|r}5 & 1 & 0 & 19 \\ 0 & 1 & 2 & 24 \\ 2 & 0 & 2 & 17 \\ 0 & 1 & 1 & 23 \\ \end{array}\right)\xrightarrow{L_{2}:\ (3),\ 2^{-1} = 14}\left(\begin{array}{rrr|r}1 & 0 & 1 & 22 \\ 0 & 1 & 2 & 24 \\ 0 & 1 & 22 & 17 \\ 0 & 1 & 1 & 23 \\ \end{array}\right)\]\[\xrightarrow{L_{3}:\ (2),\ 1^{-1} = 1}\left(\begin{array}{rrr|r}1 & 0 & 1 & 22 \\ 0 & 1 & 2 & 24 \\ 0 & 0 & 20 & 20 \\ 0 & 0 & 26 & 26 \\ \end{array}\right)\xrightarrow{L_{5}:\ (4),\ 26^{-1} = 26}\left(\begin{array}{rrr|r}1 & 0 & 0 & 21 \\ 0 & 1 & 0 & 22 \\ 0 & 0 & 1 & 1 \\ 0 & 0 & 0 & 0 \\ \end{array}\right)\]$\Rightarrow$ $L_{2} \equiv 21$, $L_{3} \equiv 22$, $L_{5} \equiv 1$ $(\textrm{mod}\ 27)$.\par\noindent Собираем по КТО (модули $4$, $27$): $L_{2} = 75$; $L_{3} = 76$; $L_{5} = 28$. Проверка: $56^{75} \equiv 2$, $56^{76} \equiv 3$, $56^{28} \equiv 5$ $(\textrm{mod}\ 109)$.\[L_{2} = 75,\ L_{3} = 76,\ L_{5} = 28.\]
\end{example}
